% appendix-numerical-v2.tex — \input'd from main.tex (after \appendix)
%
% Rewrite of Appendix A: numerical implementation, pipeline-ordered.
% Single file; replaces appendix-numerical.tex + three sub-files.
%
% Sources:
%   - Approved content from appendix-numerical-combinatorics.tex (A.1)
%   - Approved content from appendix-numerical-dismissal.tex (A.2)
%   - Rewritten A.3 (reads well but evolving)
%   - New: intro paragraph, perturbation trick, Q error bound,
%     accumulator with error bounds (all unreviewed)
%
% Structure:
%   Intro paragraph
%   B.1 Input Representation
%   B.2 Transition Feasibility and Pruning
%   B.3 KKT System and Numerical Error
%   B.4 Accumulator and Final Answer

\section{Numerical Implementation}\label{app:numerical-v2}

% --- Intro paragraph (NEW, unreviewed) ---

The algorithms in Section~\ref{sec:algorithms} are stated for exact
real arithmetic: the KKT system is solved without error, and
predicates such as $\beta_i > 0$ are decided exactly.
An exact implementation over~$\mathbb{Q}$ is possible
(Section~\ref{app:v2-input}), but prohibitively slow for the
exponential enumeration in Algorithm~\ref{alg:ehz}.
We therefore use \texttt{f64} floating-point arithmetic for the
continuous computations (eigendecomposition, KKT solve) and reserve exact
arithmetic for discrete decisions (combinatorial data,
sign predicates).

This introduces two difficulties.
First, \emph{branching decisions} based on numerical values
(e.g.\ ``is $\beta_i > 0$?'') can go either way when the true
value is near the threshold; the algorithm must handle
the resulting uncertainty.
Second, \emph{numerically unstable steps}---in particular,
the eigendecomposition solve of near-singular KKT systems---can
produce large errors in the solution vector~$\hat\beta$, though
we show that the error in the objective value~$\hat Q$ remains
controlled.
Since we seek only an approximate minimum action
$A_{\min} = 1/(2\,Q_{\max})$, controlled~$Q$ error suffices.

The remainder of this appendix discusses these issues following
the algorithm's pipeline: input representation
(Section~\ref{app:v2-input}), transition feasibility and pruning
(Section~\ref{app:v2-pruning}), the KKT solve and its numerical
error (Section~\ref{app:v2-kkt}), and the accumulator that
combines all candidates into a final answer
(Section~\ref{app:v2-answer}).


% ===================================================================
\subsection{Input Representation}\label{app:v2-input}
% ===================================================================

% --- Rational halfspace data (from approved A.1) ---

\paragraph{Rational halfspace data.}
The algorithms in Section~\ref{sec:algorithms} operate on a polytope
$K = K(n, h) = \{x \in \mathbb{R}^4 : \langle n_i, x \rangle \leq h_i,
\; i = 1, \ldots, F\}$.
Our input polytopes have rational halfspace data:
normals $n_i \in \mathbb{Q}^4 \setminus \{0\}$
(not necessarily unit vectors, unlike the unit normals
in Section~\ref{sec:algorithms})
and heights $h_i \in \mathbb{Q}_{>0}$.
This allows several combinatorial quantities to be
computed exactly over~$\mathbb{Q}$, eliminating
numerical predicates from all discrete decisions.

% --- Exact quantities (from approved A.1, simplified per Jörn) ---

\begin{remark}[Exact quantities from rational coordinates]%
\label{rem:v2-exact-quantities}
Given a polytope $K = K(n, h)$ with rational halfspace data,
the following quantities are computed exactly:
\begin{enumerate}
\item \emph{Vertices.}
  For each 4-element subset $S \subset \{1, \ldots, F\}$
  with $\det(N_S) \neq 0$ (where $N_S$ is the $4 \times 4$
  matrix with rows~$n_i^\top$, $i \in S$),
  the candidate vertex $v_S = N_S^{-1} h_S$
  (where $h_S = (h_i)_{i \in S}$) is computed
  via Cramer's rule over~$\mathbb{Q}$.
  The subset~$S$ defines a vertex of~$K$ if
  $\langle n_i, v_S \rangle \leq h_i$ for all $i \notin S$.
\item \emph{Vertex--facet incidence.}
  The equality $\langle n_i, v_j \rangle = h_i$
  is exactly decidable over~$\mathbb{Q}$.
\item \emph{Symplectic signs.}
  For each pair of facets $i, k$,
  the value $\omega_0(n_i, n_k) \in \mathbb{Q}$
  is exact (since $n_i \in \mathbb{Q}^4$),
  so $\operatorname{sign}(\omega_0(n_i, n_k))
  \in \{+, -, 0\}$ is determined without
  numerical error.
\end{enumerate}
\end{remark}

% --- Perturbation trick (NEW, unreviewed) ---

\paragraph{Generic perturbation.}
% [TODO: JÖRN - verify the continuity claim for c_EHZ on polytopes.
%  vol continuity is clear (continuous function of vertices).
%  c_EHZ continuity w.r.t. Hausdorff distance is a known result
%  but should be cited.]
We may assume without loss of generality that
$\omega_0(n_i, n_k) \neq 0$ for all pairs of facets $i \neq k$.
This is achieved by perturbing the rational normals:
replace each~$n_i$ by $n_i + \delta n_i$, where the
$\delta n_i \in \mathbb{Q}^4$ are chosen randomly with
$\|\delta n_i\| / \|n_i\| < 2^{-64}$.
Since $2^{-64}$ is well below the \texttt{f64} machine
epsilon~$\approx 2^{-52}$, the floating-point representations
of the normals and heights are unchanged:
all \texttt{f64} computations produce identical results
on the original and perturbed polytopes.

The condition $\omega_0(n_i, n_k) = 0$ defines an algebraic
hypersurface of codimension~$\geq 1$ in the space of normal
vectors.
Since this hypersurface has Lebesgue measure zero
in~$\mathbb{R}^{4F}$, there is no a~priori reason for
a pseudorandom rational perturbation to land on it;
in practice, all perturbations we generate satisfy
$\omega_0(n_i, n_k) \neq 0$ for all pairs $i \neq k$.
% [TODO: JÖRN - do we also want simplicity (each vertex on
%  exactly 4 facets) from the perturbation? If so, the
%  non-simplicity locus is also codimension >= 1. But I'm
%  unsure if simplicity is used anywhere in the algorithm.]

Since $\operatorname{vol}$ and $c_{\mathrm{EHZ}}$ are continuous
functions of the halfspace data, the perturbation changes
the systolic ratio by a negligible amount.
The sign pattern
$\operatorname{sign}(\omega_0(n_i, n_k)) \in \{+1, -1\}$
(with no zeros) is now fully determined from the rational
data, simplifying the enumeration of candidate
pairs~$(S, \sigma)$ and the transition feasibility check below.
% [TODO: JÖRN - for the billiard algorithm on Lagrangian products,
%  the theorems (billiard characterization, bounce bound) need to
%  carry over to infinitesimal perturbations of Lagrangian products.
%  This is out of scope for this appendix but should be noted.]

% --- F64 conversion (from approved A.1) ---

\begin{remark}[Floating-point conversion]%
\label{rem:v2-f64-conversion}
A rational polytope $K = K(n, h)$ is converted to
floating-point by normalizing each normal to a unit
vector:
\[
  \hat n_i = \frac{n_i}{\|n_i\|}, \qquad
  \hat h_i = \frac{h_i}{\|n_i\|}.
\]
The halfspace $\langle n_i, x \rangle \leq h_i$ is
equivalent to
$\langle \hat n_i, x \rangle \leq \hat h_i$.
Both the rational-to-floating-point conversion
and the normalization introduce rounding errors.
The capacity algorithm uses:
\begin{itemize}
\item the \emph{exact} rational data for discrete
  decisions (which facet subsets to enumerate,
  the signs of~$\omega_0(n_i, n_k)$,
  and the exact part of transition pruning),
\item the \emph{rounded} floating-point coordinates
  for continuous computation (the eigendecomposition-based
  KKT solver and the LP feasibility check below).
\end{itemize}
\end{remark}


% ===================================================================
\subsection{Transition Feasibility and Pruning}\label{app:v2-pruning}
% ===================================================================

% --- Transition feasibility lemma (from approved A.1) ---

\begin{lemma}[Transition feasibility]%
\label{lem:v2-transition-feasibility}
Let $K \subset \mathbb{R}^4$ be a convex polytope with
facets $F_1, \ldots, F_m$ and outward normals~$n_k$
(not necessarily unit).
A \emph{transition} $F_i \to F_j$ is a segment of a
Reeb trajectory that travels with velocity~$R_i$
(Definition~\ref{def:reeb-vectors}) for positive time,
then immediately with velocity~$R_j$ for positive time.
A transition $F_i \to F_j$ exists if and only if
\begin{enumerate}
\item $\omega_0(n_i, n_j) \ge 0$, and
\item $\exists\, x \in F_i \cap F_j$ with
  $\langle n_k, x \rangle < h_k$
  for all $k \notin \{i,j\}$ satisfying
  $\omega_0(n_i, n_k) < 0$ or
  $\omega_0(n_j, n_k) > 0$.
\end{enumerate}
\end{lemma}

\begin{proof}
A transition $F_i \to F_j$ exists if and only if
there is a point $x \in F_i \cap F_j$ and
$\varepsilon > 0$ such that
$x + t\, J_0 n_i \in K$ for $t \in [-\varepsilon, 0]$
and $x + t\, J_0 n_j \in K$ for $t \in [0, \varepsilon]$:
since $R_i = (2/h_i)\, J_0\, n_i$ with $h_i > 0$,
traveling with velocity~$R_i$ is equivalent to
moving in direction~$J_0 n_i$.

\emph{Necessity.}
Let $x \in F_i \cap F_j$ realize the transition.
For the approach ($t < 0$) to satisfy
the~$F_j$ constraint:
$\langle n_j, x + t\, J_0 n_i \rangle
= h_j + t\,\omega_0(n_i, n_j) \leq h_j$,
which forces $\omega_0(n_i, n_j) \geq 0$
(since $t < 0$).
For the departure ($t > 0$), the~$F_i$ constraint gives
$\langle n_i, x + t\, J_0 n_j \rangle
= h_i - t\,\omega_0(n_i, n_j) \leq h_i$,
again requiring $\omega_0(n_i, n_j) \geq 0$.
This proves~(1).

For~(2), suppose $k \notin \{i,j\}$ satisfies
$\omega_0(n_i, n_k) < 0$.
Then
$\langle n_k, x + t\, J_0 n_i \rangle
= \langle n_k, x \rangle + t\,\omega_0(n_i, n_k)$.
Since $t < 0$ and $\omega_0(n_i, n_k) < 0$,
their product $t\,\omega_0(n_i, n_k) > 0$, so
$\langle n_k, x + t\, J_0 n_i \rangle
> \langle n_k, x \rangle$.
If $\langle n_k, x \rangle = h_k$, this
exceeds~$h_k$, violating $x + t\, J_0 n_i \in K$.
Hence $\langle n_k, x \rangle < h_k$.
The case $\omega_0(n_j, n_k) > 0$ is analogous
for the departure.

\emph{Sufficiency.}
Given $x \in F_i \cap F_j$ satisfying~(1) and~(2),
we construct $\varepsilon > 0$ witnessing
the transition.
For $k \notin \{i,j\}$ with
$\omega_0(n_i, n_k) \geq 0$ and
$\omega_0(n_j, n_k) \leq 0$:
$\langle n_k, x + t\, J_0 n_i \rangle
\leq \langle n_k, x \rangle \leq h_k$
for $t \leq 0$,
and
$\langle n_k, x + t\, J_0 n_j \rangle
\leq \langle n_k, x \rangle \leq h_k$
for $t \geq 0$.
For the remaining $k \notin \{i,j\}$
(those with $\omega_0(n_i, n_k) < 0$ or
$\omega_0(n_j, n_k) > 0$),
condition~(2) gives
$\langle n_k, x \rangle < h_k$.
Let $\delta$ be the minimum of
$h_k - \langle n_k, x \rangle > 0$
over these~$k$ (or any $\delta > 0$ if there are none),
and choose
$\varepsilon < \delta / \max_k(
|\omega_0(n_i, n_k)| + |\omega_0(n_j, n_k)|)$.
Then all halfspace constraints hold for the
approach and departure, so $x$ together
with~$\varepsilon$ witnesses the transition
$F_i \to F_j$.
\end{proof}

% --- Numerical pruning (from approved A.1, reframed) ---

\paragraph{Early-exit heuristic.}
The transition feasibility lemma provides a sound
early-exit rule: if we can certify that a transition
$F_{\sigma(i)} \to F_{\sigma(i+1)}$ does not exist for
some consecutive pair in the cycle~$\sigma$, then no
closed Reeb orbit with combinatorics~$(S, \sigma)$ exists,
and the pair is skipped.

Condition~(1) is checked exactly using the rational
symplectic signs (Remark~\ref{rem:v2-exact-quantities}).
After the generic perturbation, $\omega_0(n_i, n_j) \neq 0$
for all pairs, so every edge in the adjacency graph is
unidirectional: for each pair $i \neq j$, exactly one of
the directions $F_i \to F_j$ or $F_j \to F_i$ satisfies
condition~(1).
This cuts down the number of candidate
cycles~$\sigma$ to enumerate.

Condition~(2) is an LP feasibility problem:
does there exist $x \in F_i \cap F_j$ strictly inside
the halfspaces indexed by the critical~$k$?
The numerical LP uses floating-point arithmetic and
checks feasibility with a relaxed constraint
$\langle n_k, x \rangle \leq h_k + \varepsilon_{\mathrm{LP}}$
for a small tolerance~$\varepsilon_{\mathrm{LP}} > 0$.
Since the relaxed feasible region contains the exact one,
infeasibility of the relaxed LP implies infeasibility
of the exact LP, so pruning is safe.
This makes the heuristic conservative:
it may keep pairs that a more precise check would prune,
but it never prunes a pair that admits a valid transition.
\begin{itemize}
\item Pruned transitions are correctly pruned
  (the heuristic provides a sufficient condition for infeasibility).
\item Transitions not pruned may still be infeasible in
  exact arithmetic; these pairs are processed anyway.
  The extra work is a performance cost, not a
  correctness issue, since the algorithm's correctness
  does not depend on the pruning step.
\end{itemize}


% ===================================================================
\subsection{KKT System and Numerical Error}\label{app:v2-kkt}
% ===================================================================

% --- Intro: symmetric KKT and eigendecomposition (unreviewed) ---

For each candidate pair $(S, \sigma)$ that survives
the pruning of Section~\ref{app:v2-pruning},
the algorithm solves the KKT linear system.
The standard stationarity condition
$H\beta = N\lambda + \eta\nu$
(cf.\ equation~\eqref{eq:linear-system})
uses multipliers~$\lambda \in \mathbb{R}^4$ and~$\nu \in \mathbb{R}$
that produce a non-symmetric coefficient matrix.
We define \emph{negated multipliers}
$\mu = -\lambda$ and $\xi = -\nu$, so that
the stationarity condition becomes
$H\beta + N\mu + \eta\xi = 0$
and the full system takes the symmetric saddle-point form
\begin{equation}\label{eq:v2-kkt-system}
  \underbrace{\begin{pmatrix} H & N & \eta \\[3pt]
    N^\top & 0 & 0 \\[3pt]
    \eta^\top & 0 & 0 \end{pmatrix}}_{M}
  \begin{pmatrix} \beta \\ \mu \\ \xi \end{pmatrix}
  = \begin{pmatrix} 0 \\ 0 \\ 1 \end{pmatrix},
\end{equation}
where $M$ is a real symmetric matrix of size
$(m + 5) \times (m + 5)$ with $m = |S|$.
Since~$M$ is real symmetric, it admits an eigendecomposition
$M = V \Lambda V^\top$ with real eigenvalues
$\lambda_1 \geq \cdots \geq \lambda_{m+5}$
and orthogonal eigenvectors $V = (v_1, \ldots, v_{m+5})$.
(No confusion with the original Lagrange multipliers arises,
since we use~$\mu$, $\xi$ throughout.)
The system is solved via eigendecomposition in double-precision
floating-point arithmetic.
The numerical solution is
$\hat x = (\hat\beta, \hat\mu, \hat\xi)
  = \sum_{|\lambda_j| > \tau}
    \frac{v_j^\top b}{\lambda_j}\, v_j$,
where the sum retains eigenvalues with magnitude above a
rank-detection threshold~$\tau$, and division by
small~$|\lambda_j|$ amplifies errors in the corresponding
directions.
This single factorisation simultaneously yields the
solution, the rank ($|\lambda_j| > \tau$ count),
the null space (eigenvectors for $|\lambda_j| \leq \tau$),
and the inertia (counts of positive, zero, and negative
eigenvalues; Lemma~\ref{lem:v2-kkt-inertia}).
Throughout this section, $\|\cdot\|$ denotes the Euclidean
($\ell^2$) norm for vectors and the spectral norm
(operator $2$-norm, equal to $\max_j |\lambda_j|$)
for matrices.

\begin{remark}[Eigendecomposition accuracy]%
\label{rem:v2-eigen-accuracy}
The computed eigendecomposition of a real symmetric matrix
is backward stable: the computed eigenvectors~$\hat V$ and
eigenvalues~$\hat\Lambda$ are the exact eigendecomposition
of a nearby matrix $M + E$ with
$\|E\| \leq p(m)\,\varepsilon_{\mathrm{mach}}\,\|M\|$,
where $p(m)$ is a modest polynomial in the matrix
dimension~$m + 5$
\cite[Ch.~8]{GolubVanLoan2013}.
In particular:
\begin{enumerate}
\item \emph{Eigenvalues.}
  By Weyl's theorem,
  $|\hat\lambda_j - \lambda_j| \leq \|E\|
    \leq p(m)\,\varepsilon_{\mathrm{mach}}\,\|M\|$.
  Large eigenvalues have high relative accuracy;
  small ones ($|\lambda_j| \ll \|M\|$) may have large
  relative error but are correctly identified as small.
\item \emph{Eigenvectors.}
  By the Davis--Kahan $\sin\theta$ theorem
  \cite[Ch.~8]{GolubVanLoan2013},
  the angular error of the $j$-th eigenvector satisfies
  \[
    \sin\theta(\hat v_j, v_j)
      \leq \frac{\|E\|}{\operatorname{gap}_j},
    \qquad
    \operatorname{gap}_j
      = \min_{k \neq j} |\lambda_j - \lambda_k|.
  \]
  Clustered eigenvalues yield poorly determined
  individual vectors, but their joint invariant subspace
  remains accurate.
\item \emph{Solve step.}
  For non-singular~$M$, the solution~$x$ is unique and
  the relative error satisfies
  $\|\hat x - x\| / \|x\|
    = O(\kappa(M)\,\varepsilon_{\mathrm{mach}})$,
  where $\kappa(M) = |\lambda_{\max}| / |\lambda_{\min}|$
  \cite[Ch.~5]{Higham2002}.
  \textbf{This bound does not apply to our KKT systems.}
  When $M$ is singular, $x$~is non-unique and
  $\|\hat x - x\|$ can be arbitrarily large
  (any null-space component is lost).
  When $M$ is near-singular,
  $\kappa(M) \gg 1$ and the bound is vacuous.
  We therefore analyse the error in~$Q$ directly
  (Remarks~\ref{rem:v2-null-q-constant}--\ref{rem:v2-near-null-q}),
  obtaining a tight bound
  (Lemma~\ref{lem:v2-q-interval})
  that exploits the insensitivity of~$Q$
  to near-null directions.
\end{enumerate}
For our matrices ($m + 5 \leq 21$, since $|S| \leq 16$),
the polynomial factor~$p(m)$ is negligible.
The dominant error source is the condition
number~$\kappa(M)$.
% [TODO: JÖRN - verify chapter references for Golub--Van Loan
%  eigenvalue perturbation theory (Weyl's theorem and
%  Davis--Kahan, both in Ch.~8 of 4th ed per agent's
%  general knowledge). The results are standard but the
%  exact section numbers should be checked.]
\end{remark}

% --- Inertia of the saddle-point matrix (unreviewed) ---

\begin{lemma}[Inertia of the KKT matrix]%
\label{lem:v2-kkt-inertia}
Let $M$ be the symmetric KKT matrix
of~\eqref{eq:v2-kkt-system} with
constraint matrix
$A = \bigl(\begin{smallmatrix} N^\top \\ \eta^\top
  \end{smallmatrix}\bigr) \in \mathbb{R}^{5 \times m}$
and $p = \operatorname{rank}(A) \leq \min(5, m)$.
Let $T = \ker(A) = \{\delta\beta \in \mathbb{R}^m :
  N^\top\delta\beta = 0,\; \eta^\top\delta\beta = 0\}$
be the \emph{tangent space} of the constraint surface,
with $\dim T = m - p$.
Let $H|_T = Z^\top H Z$ denote the restriction of~$H$
to~$T$, where $Z \in \mathbb{R}^{m \times (m-p)}$ is any
matrix whose columns form a basis of~$T$.

Then the eigenvalue inertia of~$M$
(the triple $(n_+, n_0, n_-)$ counting positive, zero,
and negative eigenvalues) satisfies
\begin{equation}\label{eq:v2-inertia}
  n_+(M) = n_+(H|_T) + p, \quad
  n_0(M) = n_0(H|_T) + (5 - p), \quad
  n_-(M) = n_-(H|_T) + p.
\end{equation}
In the generic case ($m \geq 5$ and $p = 5$), this
simplifies to
$n_+(M) = n_+(H|_T) + 5$ and
$n_-(M) = n_-(H|_T) + 5$.
\end{lemma}

\begin{proof}
% [TODO: JÖRN - agent verified the statement numerically
%  for all (S,σ) nodes of all test polytopes (Table 8--9
%  of q_error). Only the result is used; proof deferred.]
See \cite[Thm.~3.2]{BenziGolubLiesen2005} for a
complete proof.
The result exceeds our scope (it requires the full
spectral theory of symmetric saddle-point matrices),
and only the inertia
formula~\eqref{eq:v2-inertia} is used below.
\end{proof}

\begin{remark}[Eigenvalue signs and the objective]%
\label{rem:v2-eigenvalue-signs}
Recall that the objective
$Q(\beta) = -(1/2)\,\beta^\top H\,\beta$
has Hessian $-H$ (a constant matrix).
The second-order behaviour of~$Q$ on the constraint
surface is governed by the restricted Hessian
$(-H)|_T = -H|_T$.
By Lemma~\ref{lem:v2-kkt-inertia}
(with $p = \operatorname{rank}(A)$),
we can read off the inertia of~$H|_T$ directly from
the eigenvalues of~$M$:
\begin{itemize}
\item $n_-(M) = p$ (only ``structural''
  negatives): $H|_T$ is positive
  semi-definite ($n_-(H|_T) = 0$).
  If also $n_0(M) = 5 - p$
  (so $n_0(H|_T) = 0$),
  then $H|_T$ is positive definite,
  $(-H)|_T$ is negative definite,
  and~$Q$ is \emph{strictly concave} on~$T$:
  $\beta_0$ is a strict local maximum of~$Q$
  on the constraint surface.

\item If $n_-(M) > p$, then $H|_T$ has at least one
  negative eigenvalue.
  The corresponding eigendirection of~$M$
  (restricted to the~$\beta$-components) is a
  direction along which~$Q$ \emph{increases}
  from~$\beta_0$.

\item If $n_0(M) > 5 - p$, then $H|_T$ has a zero
  eigenvalue and~$Q$ is \emph{constant} (to first
  order) along the corresponding tangent direction.
  This direction lies in the null space of~$M$
  (Remark~\ref{rem:v2-null-q-constant}).
\end{itemize}
In the generic case ($m \geq 5$, $p = 5$, and $H|_T$
non-degenerate), the classification reduces to
counting $n_-(M)$:
\begin{gather*}
  n_-(M) = 5 \\
  \Updownarrow \\
  \text{$\beta_0$ is a strict local maximum
  of~$Q$ on the constraint surface.}
\end{gather*}
\end{remark}

\begin{remark}[Null-space directions preserve~$Q$]%
\label{rem:v2-null-q-constant}
Let $v = (v_\beta, v_\mu, v_\xi)$ be an eigenvector of~$M$
with eigenvalue~$0$, so $Mv = 0$.
Then $N^\top v_\beta = 0$ and $\eta^\top v_\beta = 0$
(from the constraint rows of~$M$),
so $v_\beta \in T$.
Moreover, $H v_\beta + N v_\mu + \eta v_\xi = 0$
(stationarity row), so
$v_\beta^\top H v_\beta
  = -v_\beta^\top(N v_\mu + \eta v_\xi)
  = -(N^\top v_\beta)^\top v_\mu
    - (\eta^\top v_\beta)\,v_\xi = 0$.
For the cross term, the KKT stationarity condition
$H\beta_0 + N\mu_0 + \eta\xi_0 = 0$ gives
$v_\beta^\top H\,\beta_0
  = -v_\beta^\top(N\mu_0 + \eta\xi_0)
  = -(N^\top v_\beta)^\top\mu_0
    - (\eta^\top v_\beta)\,\xi_0 = 0$
(since $v_\beta \in T$).
Hence $v_\beta^\top H\,\beta_0 = 0$
and $v_\beta^\top H\,v_\beta = 0$, so
\[
  Q(\beta_0 + t\,v_\beta) - Q(\beta_0)
  = -t\,v_\beta^\top H\,\beta_0
    - \tfrac{1}{2}\,t^2\,v_\beta^\top H\,v_\beta
  = 0
\]
for all~$t$.
In other words, $Q$ is exactly constant along
null-space directions.
The same reasoning shows that $Q$ is
constraint-preserving:
the displaced~$\beta_0 + t\,v_\beta$ satisfies
$N^\top(\beta_0 + t\,v_\beta) = 0$ and
$\eta^\top(\beta_0 + t\,v_\beta) = 1$.
\end{remark}

\begin{remark}[$Q$ variation along near-null directions]%
\label{rem:v2-near-null-q}
For an eigenvalue~$\lambda_j$ of~$M$ that is
small but nonzero, the eigenvector
$v_j = (v_\beta, v_\mu, v_\xi)$ satisfies
$N^\top v_\beta = \lambda_j\,v_\mu$ and
$\eta^\top v_\beta = \lambda_j\,v_\xi$,
so $v_\beta$ is \emph{approximately} tangent
(constraint violation of order~$|\lambda_j|$).
The~$Q$ variation along~$v_\beta$ is
\begin{align*}
  Q(\beta_0 + t\,v_\beta) - Q(\beta_0)
  &= -t\,v_\beta^\top H\,\beta_0
    - \tfrac{1}{2}\,t^2\,v_\beta^\top H\,v_\beta.
\end{align*}
From the eigenvector equation
$Hv_\beta = \lambda_j\,v_\beta
  - N v_\mu - \eta v_\xi$,
one computes
\begin{equation}\label{eq:v2-near-null-curvature}
  v_\beta^\top H\,v_\beta
  = \lambda_j\bigl(2\,\|v_\beta\|^2 - 1\bigr),
\end{equation}
using the orthogonality
$\|v_\beta\|^2 + \|v_\mu\|^2 + v_\xi^2 = 1$
(the eigenvector is unit-length).
Both the linear and quadratic terms are
proportional to~$\lambda_j$:
\begin{itemize}
\item The linear term:
  the KKT condition
  $H\beta_0 = -N\mu_0 - \eta\xi_0$ gives
  $v_\beta^\top H\,\beta_0
    = -(N^\top v_\beta)^\top\mu_0
      - (\eta^\top v_\beta)\,\xi_0
    = -\lambda_j(v_\mu^\top\mu_0
      + v_\xi\,\xi_0)$,
  which is $O(|\lambda_j|)$.
\item The curvature~\eqref{eq:v2-near-null-curvature}
  is $O(|\lambda_j|)$.
\end{itemize}
In particular, directions with small~$|\lambda_j|$
change~$Q$ only slowly.
This is the mechanism underlying the $Q$~error bound
(Lemma~\ref{lem:v2-q-interval}): even though the
numerical error~$\delta\beta$ can be large along
near-null directions (of order~$\|r\| / |\lambda_j|$),
the resulting~$Q$ error is only second-order in~$\|r\|$
because the $Q$~gradient along these directions
is itself $O(|\lambda_j|)$,
partially cancelling the $1/|\lambda_j|$ amplification.
\end{remark}

\begin{remark}[Near-singularity and error decomposition]%
\label{rem:v2-error-decomposition}
By the generic perturbation (Section~\ref{app:v2-input}),
the exact KKT system is generically non-singular.
However, it can be \emph{near}-singular:
the smallest eigenvalue magnitude
$|\lambda_{\min}| = \min_j |\lambda_j|$ may be close
to zero, causing large numerical error in the
solution~$(\hat\beta, \hat\mu, \hat\xi)$.

The numerical error in~$\hat\beta$ decomposes along the
eigenvectors of~$M$ into two qualitatively different
contributions:
\begin{itemize}
\item \emph{Near-null directions}
  (small~$|\lambda_j|$):
  The error in~$\hat x$ can be large
  (absolute error~$\sim |\lambda_{\max}|/|\lambda_{\min}|$;
  Remark~\ref{rem:v2-eigen-accuracy}~(3)).
  However, the~$Q$ error is only
  \emph{second-order} in the perturbation
  (Remark~\ref{rem:v2-near-null-q}):
  the near-null eigenvector is approximately tangent
  and has curvature proportional to~$\lambda_j$,
  so the large displacement~$\sim 1/|\lambda_j|$
  is partially cancelled by the
  small curvature~$\sim |\lambda_j|$.
  The sign of~$\lambda_j$ determines whether
  this direction is concave or convex for~$Q$
  (Remark~\ref{rem:v2-eigenvalue-signs}).

\item \emph{Well-conditioned directions}
  (large~$|\lambda_j|$):
  The solve is accurate
  ($\|\delta\beta\| = O(\varepsilon_{\mathrm{mach}})$),
  so
  $|Q(\hat\beta) - Q(\beta_0)|
    = O(\varepsilon_{\mathrm{mach}})$.
\end{itemize}
\end{remark}

% --- Q error bound lemma (replaces dismiss logic B.5--B.9, unreviewed) ---

\begin{lemma}[$Q$ error bound from KKT residual]%
\label{lem:v2-q-interval}
Let $M$ be the symmetric KKT matrix
of~\eqref{eq:v2-kkt-system}
with smallest eigenvalue magnitude
$|\lambda_{\min}| = \min_j |\lambda_j| > 0$,
and let $x_0 = (\beta_0, \mu_0, \xi_0)$ be the exact
solution (with negated multipliers
$\mu_0 = -\lambda_0$, $\xi_0 = -\nu_0$).
Let $\hat x = (\hat\beta, \hat\mu, \hat\xi)$ be the
numerical solution with residual
\[
  r \;=\; M\hat x - b
  \;=\; \begin{pmatrix}
    H\hat\beta + N\hat\mu + \eta\hat\xi \\
    N^\top\hat\beta \\
    \eta^\top\hat\beta - 1
  \end{pmatrix}
  \;=:\; \begin{pmatrix}r_1 \\ r_2 \\ r_3\end{pmatrix},
\]
where $r_1 \in \mathbb{R}^m$ is the stationarity residual,
$r_2 \in \mathbb{R}^4$ is the constraint residual, and
$r_3 \in \mathbb{R}$ is the normalisation residual.
Define the \emph{residual-corrected objective}
\begin{equation}\label{eq:v2-q-corrected}
  \tilde Q
  \;=\; Q(\hat\beta)
    - \bigl(r_2^\top\hat\mu + r_3\,\hat\xi\bigr)
\end{equation}
and the \emph{second-order error bound}
\begin{equation}\label{eq:v2-q-error-bound}
  E \;=\; \frac{\|r\|^2}{|\lambda_{\min}|}
    \Bigl(2 + \frac{\|H\|}{2\,|\lambda_{\min}|}\Bigr).
\end{equation}
Then $|Q(\beta_0) - \tilde Q| \leq E$.
\end{lemma}

\begin{proof}
We show $Q(\beta_0) = \tilde Q + R$ with $|R| \leq E$
by expanding the difference $Q(\hat\beta) - Q(\beta_0)$,
using the KKT condition to rewrite the cross term,
then substituting numerical for exact multipliers.

Write $\delta x = \hat x - x_0$ with components
$(\delta\beta, \delta\mu, \delta\xi)$.
Since $M x_0 = b$, we have $M\,\delta x = r$,
hence $\|\delta x\| \leq \|r\| / |\lambda_{\min}|$ and
each component satisfies the same bound:
\begin{equation}\label{eq:v2-delta-bound}
  \|\delta\beta\| \leq \frac{\|r\|}{|\lambda_{\min}|}, \quad
  \|\delta\mu\| \leq \frac{\|r\|}{|\lambda_{\min}|}, \quad
  |\delta\xi| \leq \frac{\|r\|}{|\lambda_{\min}|}.
\end{equation}

\emph{Step~1: expand $Q(\hat\beta) - Q(\beta_0)$.}
Since $Q(\beta) = -(1/2)\,\beta^\top H\,\beta$,
\[
  Q(\hat\beta) - Q(\beta_0)
  = -\delta\beta^\top H\,\beta_0
    - \tfrac{1}{2}\,\delta\beta^\top H\,\delta\beta.
\]

\emph{Step~2: rewrite the cross term using KKT.}
The stationarity condition (first block of $M x_0 = b$) is
$H\beta_0 + N\mu_0 + \eta\xi_0 = 0$
(equivalently, $H\beta_0 = N\lambda_0 + \eta\nu_0$
in the original multipliers).
Hence
\[
  \delta\beta^\top H\,\beta_0
  = -\delta\beta^\top(N\mu_0 + \eta\xi_0)
  = -(N^\top\delta\beta)^\top \mu_0
    - (\eta^\top\delta\beta)\,\xi_0.
\]
The constraint residual blocks
(rows $m{+}1$ to $m{+}5$ of $M\hat x - b$) give
$N^\top\delta\beta = N^\top\hat\beta - N^\top\beta_0
  = N^\top\hat\beta = r_2$
(since $N^\top\beta_0 = 0$ by the exact constraint)
and similarly $\eta^\top\delta\beta = r_3$, hence
\begin{equation}\label{eq:v2-cross-term}
  -\delta\beta^\top H\,\beta_0
  = r_2^\top\mu_0 + r_3\,\xi_0.
\end{equation}
Substituting~\eqref{eq:v2-cross-term} into Step~1
and rearranging gives
$Q(\beta_0) = Q(\hat\beta) - (r_2^\top\mu_0 + r_3\,\xi_0)
  + \tfrac{1}{2}\,\delta\beta^\top H\,\delta\beta$.

\emph{Step~3: substitute numerical multipliers.}
Write $\mu_0 = \hat\mu - \delta\mu$ and
$\xi_0 = \hat\xi - \delta\xi$:
\begin{align*}
  r_2^\top\mu_0 + r_3\,\xi_0
  &= (r_2^\top\hat\mu + r_3\,\hat\xi)
    - (r_2^\top\delta\mu + r_3\,\delta\xi).
\end{align*}
The first term is the computable correction
in~\eqref{eq:v2-q-corrected}.
The second is a \emph{second-order remainder}:
\[
  |r_2^\top\delta\mu + r_3\,\delta\xi|
  \leq \|r_2\|\,\|\delta\mu\| + |r_3|\,|\delta\xi|
  \leq 2\,\frac{\|r\|^2}{|\lambda_{\min}|},
\]
using $\|r_2\|, |r_3| \leq \|r\|$
and~\eqref{eq:v2-delta-bound}.

\emph{Step~4: bound the quadratic term.}
\[
  \tfrac{1}{2}\,|\delta\beta^\top H\,\delta\beta|
  \leq \tfrac{1}{2}\,\|H\|\,\|\delta\beta\|^2
  \leq \frac{\|H\|\,\|r\|^2}{2\,|\lambda_{\min}|^2}.
\]

\emph{Conclusion.}
Combining Steps~1--4:
\[
  Q(\beta_0)
  = Q(\hat\beta)
    - (r_2^\top\hat\mu + r_3\,\hat\xi)
    \;+\; R,
\]
where $|R| \leq 2\,\|r\|^2 / |\lambda_{\min}|
  + \|H\|\,\|r\|^2 / (2\,|\lambda_{\min}|^2) = E$.
\qedhere
\end{proof}

\begin{algorithm}[$Q$ value with error bound]%
\label{alg:v2-q-interval}
For each pair $(S, \sigma)$ that survives pruning,
the algorithm computes a corrected objective
value~$\tilde Q$ with error bound~$E$
satisfying $|Q(\beta_0) - \tilde Q| \leq E$:
\begin{enumerate}
\item Compute the eigendecomposition
  $M = V \Lambda V^\top$ of the symmetric KKT
  matrix~\eqref{eq:v2-kkt-system}.
\item Compute the pseudoinverse solution
  $\hat x = (\hat\beta, \hat\mu, \hat\xi)
    = \sum_{|\lambda_j| > \tau}
      (v_j^\top b / \lambda_j)\, v_j$,
  retaining eigenvalues above a threshold~$\tau$.
\item Compute the residual $r = M\hat x - b$ and
  extract blocks $r_2 = N^\top\hat\beta$
  (constraint residual)
  and $r_3 = \eta^\top\hat\beta - 1$
  (normalisation residual).
\item Compute $Q(\hat\beta)$ and the corrected value
  $\tilde Q = Q(\hat\beta)
    - (r_2^\top\hat\mu + r_3\,\hat\xi)$.
\item Compute the error bound
  $E = \|r\|^2\,(2/|\lambda_{\min}|
    + \|H\| / (2\,|\lambda_{\min}|^2))$,
  where $|\lambda_{\min}|$ is the smallest
  retained eigenvalue magnitude.
\item Return the corrected value~$\tilde Q$,
  the error bound~$E$, and~$\hat\beta$.
\end{enumerate}
\end{algorithm}

\begin{remark}[Error bound in practice]%
\label{rem:v2-interval-width}
If $|\lambda_{\min}|$ is bounded away from zero,
the residual is
$\|r\| = O(\varepsilon_{\mathrm{mach}}\,\|M\|)$
(Remark~\ref{rem:v2-eigen-accuracy}), so
$E = O(\varepsilon_{\mathrm{mach}}^2)$ and
$\tilde Q \approx Q(\hat\beta)$.
The error bound is tight and the correction negligible.

For near-singular systems
($|\lambda_{\min}| \ll 1$), the residual remains small
(backward stability of the eigendecomposition), but
$E$ grows as $1/|\lambda_{\min}|^2$.
The correction~$r_2^\top\hat\mu + r_3\,\hat\xi$
absorbs the dominant first-order error, leaving only
a second-order remainder.
Without this correction, a naive bound
$|Q(\hat\beta) - Q(\beta_0)| \leq \Delta$ would require
$\Delta = O(\|r\| / |\lambda_{\min}|)$ (first-order),
which is much larger.
\end{remark}

% --- Positivity check (from approved A.2, ad-hoc framing) ---

\paragraph{Admissibility check.}
In addition to the $Q$~value with error bound,
Algorithm~\ref{alg:ehz} requires
$\beta_i > 0$ for all $i \in S$.
Since~$\hat\beta$ differs from the true~$\beta^*$ by
numerical error, a component~$\hat\beta_i$ near zero is
ambiguous: the true value might be positive, negative,
or exactly zero.
The implementation therefore returns a three-valued verdict
for each component:
\[
  \tilde P_{>0}(\hat\beta_i) =
  \begin{cases}
    \textsc{true}
      & \text{if } \hat\beta_i > \varepsilon_\beta, \\
    \textsc{false}
      & \text{if } \hat\beta_i < -\varepsilon_\beta, \\
    \textsc{unknown}
      & \text{if } |\hat\beta_i| \leq \varepsilon_\beta,
  \end{cases}
\]
where $\varepsilon_\beta$ bounds the solve error.
If all components return $\textsc{true}$, the candidate
is certified admissible.
If any component returns $\textsc{false}$, the candidate
is inadmissible and discarded.
Otherwise, some component is $\textsc{unknown}$ and
the candidate's admissibility is uncertain:
\begin{itemize}
\item If $0 < \beta^*_i < \varepsilon_\beta$:
  the pair is genuinely admissible but may be incorrectly
  discarded ($\hat\beta_i$ could round negative).
\item If $-\varepsilon_\beta < \beta^*_i < 0$:
  the pair is genuinely inadmissible but may be incorrectly
  kept ($\hat\beta_i$ could round positive).
\item If $\beta^*_i = 0$:
  the pair is covered by a smaller
  pair $(S \setminus \{i\}, \sigma|_{S \setminus \{i\}})$,
  which the algorithm also examines.
\end{itemize}

% [TODO: JÖRN - in the new framework, the three-valued
%  verdict feeds into the Q error bound: TRUE admissibility
%  gives a tight bound from numerical error alone. UNKNOWN
%  admissibility gives a wider bound or flags the node for
%  special handling. The exact interaction needs to be
%  worked out.]


% ===================================================================
\subsection{Accumulator and Final Answer}\label{app:v2-answer}
% ===================================================================

% --- Accumulator (rewritten to match code, unreviewed) ---

The algorithm processes all candidate pairs
$(S, \sigma)$ and selects the one with smallest
action (equivalently, largest~$\tilde Q$).
The accumulator tracks two candidates:
\begin{itemize}
\item \textbf{Best certified:}
  the candidate with smallest action among those whose
  numerical~$\hat\beta$ satisfies
  $\hat\beta_i > \varepsilon_\beta$ for all~$i \in S$
  (all admissibility predicates return $\textsc{true}$).
\item \textbf{Best uncertain:}
  the candidate with smallest action among those whose
  numerical~$\hat\beta$ satisfies
  $\hat\beta_i > -\varepsilon_\beta$ for all~$i \in S$
  (all predicates return $\textsc{true}$ or
  $\textsc{unknown}$, i.e.\ no predicate returns
  $\textsc{false}$).
\end{itemize}
Every certified candidate is also uncertain, so
the best uncertain action is always~$\leq$
the best certified action.

\paragraph{Candidate classification.}
Each pair $(S, \sigma)$ is classified as follows:
\begin{description}
\item[Skipped.]
  Transition pruning
  (Section~\ref{app:v2-pruning})
  determines that no Reeb trajectory with
  combinatorics~$(S, \sigma)$ exists.

\item[Inadmissible.]
  The KKT solve returns no solution
  (no $\hat\beta$ with
  $\tilde Q > \varepsilon_Q$), or some
  $\hat\beta_i < -\varepsilon_\beta$
  (a predicate returns $\textsc{false}$).
  The pair is discarded.

\item[Certified.]
  All $\hat\beta_i > \varepsilon_\beta$.
  The pair competes for best certified
  and best uncertain.

\item[Uncertain.]
  All $\hat\beta_i > -\varepsilon_\beta$
  but some $|\hat\beta_i| \leq \varepsilon_\beta$.
  The pair competes only for best uncertain.
\end{description}

\paragraph{Comparison.}
Candidates are compared by action
$A = 1/(2\,\tilde Q)$,
using point comparison on the
residual-corrected~$\tilde Q$.
% Downstream: action = 0.5 / result.q_corrected (lower action = higher Q = better orbit)
% [TODO: JÖRN - the E < 10^{-13} claim below is from
%  the diagnostic agent's run on all fixture polytopes.
%  It will be verified via q_error experiment output once
%  the full experiment bundle is wired up.]
The error bound~$E$ from
Lemma~\ref{lem:v2-q-interval}
is tracked but not used for interval majorization:
in all test polytopes (33~entries covering
5--8~facets), the bound satisfies
$E < 10^{-13}$, so point comparison on~$\tilde Q$
is reliable.

\paragraph{Final answer.}
The algorithm returns the best certified candidate.
If no certified candidate exists, the algorithm
returns~\textsc{None} (no admissible closed Reeb
orbit was found).

\paragraph{Safety net.}
After enumeration, the algorithm checks that the best
uncertain action is~$\geq$ the best certified action.
If an uncertain candidate achieves strictly lower action
than the best certified one, the reported capacity might
be wrong---an orbit that is genuinely admissible
(with some~$\beta_i$ very close to zero) would give
a larger capacity than the certified answer.
The algorithm panics in this case rather than returning
a potentially incorrect result.
% Downstream: assert!(uncertain_cap >= certified.0, ...)
% in hk2017/mod.rs:146-152
In all experiments to date (33~test polytopes),
the assertion has never fired.

\begin{remark}[Why point comparison suffices]%
\label{rem:v2-point-comparison}
The corrected~$\tilde Q$ absorbs the first-order
numerical error via the residual correction
(Lemma~\ref{lem:v2-q-interval}).
The remaining error~$E$ is second-order
in the backward error and satisfies
$E = O(\varepsilon_{\mathrm{mach}}^2)$ for
well-conditioned systems
(Remark~\ref{rem:v2-interval-width}).
For the polytopes in our dataset,
$E < 10^{-13}$ in all cases, making interval
majorization unnecessary: point comparison
on~$\tilde Q$ correctly identifies the global
maximum.
If future polytopes produce larger~$E$ (e.g.\ due
to near-singular KKT systems),
the interval framework from
Section~\ref{app:v2-kkt} provides a fallback:
compare $[\tilde Q - E,\; \tilde Q + E]$ intervals
and resolve overlaps.
\end{remark}

\begin{remark}[Parameter scales in practice]\label{rem:v2-parameter-scales}
For polytopes generated by our random sampling
(Section~\ref{sec:random-sweep}),
the relevant quantities have the following typical magnitudes:
the symplectic inner products
$\omega_0(n_i, n_j) \in [-1, 1]$ (unit normals),
with $\omega_0(n_i, n_j) \approx 0$ occurring primarily
for near-Lagrangian 2-faces;
heights $h_i \in [h_{\min}, h_{\max}]$ with typical range
$[0.1, 10]$;
vertex norms $\|v\| \lesssim 100$;
and empirically $\beta_i = O(\|v\|)$, though the precise
scaling of $\beta$ with the polytope geometry (via the
constraints $N^\top \beta = 0$, $\eta^\top \beta = 1$)
deserves further investigation.
Before the normalisation $\eta^\top \beta = 1$,
the un-normalised period~$T$ of the candidate trajectory is bounded below
by the inradius (determined by $\min_k h_k$) and above by
the circumradius (determined by $\max_k \|v_k\|$);
the normalisation rescales $\beta$ by $1/T$.

These scales determine how aggressively the three-valued
predicates can distinguish $\textsc{true}$ from
$\textsc{unknown}$:
tolerances that are small relative to these scales yield
few $\textsc{unknown}$ verdicts.
\end{remark}
